\documentclass[11pt,oneside]{scrbook}
% napkin-style.sty detects the class automatically.

\usepackage{napkin-style}
\usepackage{napkin-macros}


\title{Title of Your Notes}
\author{Your Name}
\date{\today}

\begin{document}
\maketitle
\tableofcontents

% ===============================================================
\chapter{Theorem Environments}
% ===============================================================

% -------------------------------------------------------
\section{Main Environments}
% -------------------------------------------------------

\prototype{The integers $\ZZ$ with addition.}

\begin{definition}
    A \vocab{group} is a set $G$ together with a binary operation $\cdot$
    satisfying associativity, identity, and inverses.
\end{definition}

\begin{theorem}
    Every subgroup of a cyclic group is cyclic.
\end{theorem}
\begin{proof}
    Let $G = \langle g \rangle$ and $H \le G$. Take the smallest positive
    $n$ with $g^n \in H$; then $H = \langle g^n \rangle$.
\end{proof}

\begin{lemma}
    The intersection of two subgroups is a subgroup.
\end{lemma}

\begin{corollary}
    The intersection of any collection of subgroups is a subgroup.
\end{corollary}

\begin{proposition}
    Every group of prime order is cyclic.
\end{proposition}

% -------------------------------------------------------
\section{Remarks and Examples}
% -------------------------------------------------------

\subsection{Examples}

\begin{example}
    The even integers $2\ZZ$ form a subgroup of $(\ZZ, +)$.
\end{example}

\begin{remark}
    Note that $\QQ$ under multiplication is \emph{not} a group because
    $0$ has no inverse.
\end{remark}

\begin{moral}
    A group is the minimal algebraic structure in which ``symmetry'' makes sense.
\end{moral}

\subsection{Other Environments}

\begin{claim}
    The center $Z(G)$ is always a normal subgroup of $G$.
\end{claim}
\begin{subproof}
    For any $z \in Z(G)$ and $g \in G$, we have $gzg\inv = gg\inv z = z \in Z(G)$.
\end{subproof}

\begin{fact}
    Every finite group of odd order is solvable (Feit--Thompson).
\end{fact}

\subsubsection{Numbered vs unnumbered}

Use \texttt{theorem} for numbered environments and \texttt{theorem*} for
unnumbered ones (not defined in Napkin by default; add manually if needed).

% -------------------------------------------------------
\section{Exercises}
% -------------------------------------------------------

\begin{exercise}
    Show that $\Zm{p}$ is a group under multiplication for any prime $p$.
\end{exercise}

\begin{soln}
    Since $p$ is prime, every nonzero element of $\ZZ/p\ZZ$ has a multiplicative
    inverse by Bezout's theorem.
\end{soln}

% ===============================================================
\chapter{Diagrams and Macros}
% ===============================================================

% -------------------------------------------------------
\section{Commutative Diagrams}
% -------------------------------------------------------

\begin{example}
    A commutative square:
\[
\begin{tikzcd}
    A \arrow[r, "f"] \arrow[d, "g"'] & B \arrow[d, "h"] \\
    C \arrow[r, "k"']                & D
\end{tikzcd}
\]
\end{example}

% -------------------------------------------------------
\section{Useful Math Macros}
% -------------------------------------------------------

\subsection{Arrows}

\begin{itemize}
    \item Injection: $f \colon A \injto B$  \quad(\texttt{\textbackslash injto})
    \item Surjection: $f \colon A \surjto B$ \quad(\texttt{\textbackslash surjto})
    \item Labeled: $A \taking{f} B$          \quad(\texttt{\textbackslash taking\{f\}})
\end{itemize}

\subsection{Shorthands}

\begin{itemize}
    \item $\ol{AB}$, $\wt{X}$, $\wh{f}$ \quad(\texttt{\textbackslash ol}, \texttt{\textbackslash wt}, \texttt{\textbackslash wh})
    \item $x \defeq y$                   \quad(\texttt{\textbackslash defeq})
    \item $\eps$, $\inv$                 \quad(\texttt{\textbackslash eps}, \texttt{\textbackslash inv})
    \item $\floor{x}$, $\ceiling{x}$, $\abs{x}$, $\norm{x}$
\end{itemize}

\subsection{Number Fields}

$\NN,\ \ZZ,\ \QQ,\ \RR,\ \CC,\ \FF,\ \EE$

\end{document}
